% Main thesis entry point.
% Based on the ECUST master thesis LaTeX template.

\documentclass[UTF-8,a4paper,twoside]{article}

\usepackage[UTF8,heading=true]{ctex}
\usepackage{setspace}
\usepackage{fancyhdr}
\usepackage{enumitem}
\setlist[enumerate]{label=(\arabic*),fullwidth,itemindent=\parindent,listparindent=\parindent,itemsep=0ex,partopsep=0pt,parsep=0ex,nosep}
\raggedbottom
\usepackage{layout}
\usepackage{geometry}
\usepackage{graphicx}
\graphicspath{{../figures/}{template/}}
\usepackage{subfig}
\makeatletter
\let\c@lofdepth\relax
\let\c@lotdepth\relax
\makeatother
\usepackage{listings}
\usepackage{float}
\usepackage[backend=biber,style=gb7714-2015,gbalign=left,gbnamefmt=lowercase]{biblatex}
\addbibresource{../docs/references.bib}
\renewcommand{\bibfont}{\zihao{-4}}
\setlength{\bibitemsep}{0pt}
\usepackage{mfirstuc}
\gMFUnocap{in}
\gMFUnocap{of}
\gMFUnocap{and}
\gMFUnocap{for}
\gMFUnocap{the}
\gMFUnocap{th}
\gMFUnocap{on}
\DeclareFieldFormat[article]{journaltitle}{\capitalisewords{#1}}
\DeclareFieldFormat[inproceedings]{booktitle}{\capitalisewords{#1}}

\AtBeginDocument{
\setlength{\abovedisplayskip}{1ex}
\setlength{\belowdisplayskip}{1ex}
}

\usepackage{fontspec}
\usepackage{caption}
\captionsetup[figure]{belowskip=-2ex}
\usepackage{bicaption}
\usepackage{amsmath,lmodern}
\usepackage{xeCJK}
\usepackage{titletoc}
\usepackage{tocloft}
\usepackage{bm}
\usepackage{booktabs}

\newcommand{\topline}{\toprule[1.5pt]}
\newcommand{\bottomline}{\bottomrule[1.5pt]}
\renewcommand{\contentsname}{\zihao{4}\bfseries{目录}}
\setlength{\cftsecindent}{0em}
\setlength{\cftsubsecindent}{0em}
\setlength{\cftsubsubsecindent}{0em}
\setlength{\cftbeforetoctitleskip}{1em}
\setlength{\cftaftertoctitleskip}{2em}
\renewcommand{\cftdotsep}{0.5}
\renewcommand{\cftsecfont}{\zihao{4}}
\renewcommand{\cftsecpagefont}{\zihao{4}}
\renewcommand{\cftsubsecfont}{\zihao{-4}}
\renewcommand{\cftsubsubsecfont}{\zihao{-4}}
\renewcommand{\cftsecleader}{\cftdotfill{\cftdotsep}}
\setlength{\cftbeforesecskip}{0em}
\setlength{\cftbeforesubsecskip}{0em}
\setlength{\cftbeforesubsubsecskip}{0em}

\numberwithin{equation}{section}
\numberwithin{table}{section}
\numberwithin{figure}{section}
\renewcommand{\theequation}{\thesection-\arabic{equation}}
\DeclareCaptionLabelSeparator{colon}{ \ \ }
\DeclareCaptionFont{chfont}{\bfseries\zihao{5}}
\DeclareCaptionFont{enfont}{\zihao{5}}
\captionsetup[table]{font=chfont}
\captionsetup[table][bi-second]{name=Table, font=enfont}
\captionsetup[figure]{font=chfont}
\captionsetup[figure][bi-second]{name=Fig., font=enfont}
\DeclareCaptionFont{subfont}{\zihao{5}}
\captionsetup[subfloat]{font=subfont}

\geometry{left=2.5cm,right=2.5cm,top=2.8cm,bottom=2.0cm,headheight=2.0cm,footskip=1.2cm}
\setmainfont{Times New Roman}
\setCJKmainfont[Path=./template/SimSun/, Extension=.TTC, AutoFakeBold=true]{SIMSUN}
\setCJKfamilyfont{zhsong}[Path=./template/SimSun/, Extension=.TTC, AutoFakeBold=true]{SIMSUN}
\setCJKfamilyfont{zhhei}[Path=./template/SimHei/, Extension=.TTF, AutoFakeBold=true]{SIMHEI}
\renewcommand{\songti}{\CJKfamily{zhsong}}
\renewcommand{\heiti}{\CJKfamily{zhhei}}
\let\kaishu\relax
\newCJKfontfamily\kaishu[Path=./template/KaiTi/, Extension=.TTF, AutoFakeBold=true]{SIMKAI}

\ctexset{
	section={
	name={第,章},
	format=\heiti\zihao{4}\centering,
	beforeskip=3ex,
	afterskip=3ex
	},
	subsection={
	format=\bfseries\zihao{-4},
	beforeskip=0.6em,
	afterskip=0.6em
	},
	subsubsection={
	format=\zihao{-4},
	beforeskip=0ex,
	afterskip=0ex
	}
}

\usepackage{parskip}
\setlength{\parskip}{0pt}
\setlength{\parindent}{2em}

\fancypagestyle{main}{
\fancyhf{}
\linespread{1.4375}
\fancyhead[LO,RE]{\kaishu\bfseries\zihao{4}{华东理工大学}\normalfont\songti\zihao{-4}{硕士学位论文}}
\fancyhead[RO,LE]{\zihao{-4}{第}\thepage{页}}
}

\fancypagestyle{pre}{
\fancyhf{}
\linespread{1.4375}
\fancyhead[LO,RE]{\kaishu\bfseries\zihao{4}{华东理工大学}\normalfont\songti\zihao{-4}{硕士学位论文}}
\fancyhead[RO,LE]{\zihao{-4}{第}\Roman{page}{页}}
}

\begin{document}

\pagestyle{pre}
\thispagestyle{pre}
\begin{center}
	\heiti\zihao{4}你的中文论文题目
\end{center}
\vspace{\baselineskip}
\begin{center}
	\zihao{4}\bfseries{摘要}
\end{center}

\zihao{-4}
这里填写中文摘要。摘要应概括研究背景、研究问题、研究方法、主要结果和结论。

\noindent\zihao{-4}{\bfseries 关键词：}关键词一；关键词二；关键词三

\newpage
\begin{center}
	\zihao{4}\bfseries Your English Thesis Title
\end{center}
\vspace{\baselineskip}
\begin{center}
	\zihao{4}\bfseries{Abstract}
\end{center}

\zihao{-4}
Write the English abstract here.

\noindent\zihao{-4}{\bfseries Keywords \ \ }keyword one; keyword two; keyword three

\newpage
\thispagestyle{pre}
\tocloftpagestyle{pre}
\begin{center}
	\tableofcontents
\end{center}
\thispagestyle{pre}

\newpage
\pagestyle{main}
\setcounter{page}{1}

\section{绪论}

\subsection{研究背景}

\subsection{研究意义}

\subsection{国内外研究现状}

\subsection{研究内容与方法}

\subsection{论文结构}

\newpage
\section{相关理论与文献综述}

\subsection{核心概念界定}

\subsection{理论基础}

\subsection{国内外研究进展}

\subsection{研究述评}

\newpage
\section{研究设计}

% 本章通常 3000–5000 字，详细说明研究方案。
% 写清楚"用什么方法、怎么做、为什么这样做"。

\subsection{研究框架}

% 用文字或流程图描述整体研究框架。
% 如果有技术路线图，可以放在 figures/ 目录中引用：
% \begin{figure}[htbp]
%   \centering
%   \includegraphics[width=0.8\textwidth]{research-framework.png}
%   \bicaption{研究框架图}{Research Framework}
%   \label{fig:framework}
% \end{figure}

\subsection{数据来源}

% 说明数据的来源、采集方式、样本量和时间范围。
% 如果使用公开数据集，注明数据集名称和获取途径。
% 如果是问卷调查，说明问卷设计、发放和回收情况。

\subsection{研究方法}

% 详细介绍所用的分析方法或模型。
% 包括方法的基本原理、选择理由、参数设置等。
% 如果涉及数学公式，使用 equation 环境：
% \begin{equation}
%   y = \beta_0 + \beta_1 x_1 + \varepsilon
%   \label{eq:model}
% \end{equation}

\subsection{实施流程}

% 按时间或步骤说明研究的实施过程。
% 可以用 itemize 列出主要步骤。

\newpage
\section{实验、分析与结果}

% 本章是论文的核心，通常 5000–8000 字。
% 先展示结果，再进行分析和讨论。

\subsection{实验设计}

% 说明实验的具体设置、变量控制和实施过程。
% 如果是问卷研究，说明量表来源和信效度检验结果。

\subsection{结果展示}

% 用图表呈现主要结果。
% 图表标题使用中英双语，引用时使用 \ref{}：
% 如表~\ref{tab:example} 所示……
%
% 表格示例：
% \begin{table}[htbp]
%   \bicaption{描述性统计结果}{Descriptive Statistics}
%   \label{tab:example}
%   \centering
%   \begin{tabular}{lcc}
%     \topline
%     变量 & 均值 & 标准差 \\
%     \midrule
%     X1   & 3.45 & 0.82 \\
%     X2   & 4.12 & 0.67 \\
%     \bottomline
%   \end{tabular}
% \end{table}

\subsection{结果分析}

% 对结果进行解释，与文献中的发现进行对比。
% 指出哪些假设得到支持，哪些没有。

\subsection{讨论}

% 讨论结果的理论和实践意义。
% 分析可能的原因，承认局限性。
% 与第 2 章的文献综述形成呼应。

\newpage
\section{总结与展望}

% 本章通常 1500–2500 字，总结全文并展望未来。

\subsection{研究结论}

% 用 3–5 条要点概括全文的核心发现。
% 每条结论对应第 4 章的一个主要结果。
% 不要引入新的数据或引用。

\subsection{研究贡献}

% 说明本研究在理论和实践两个层面的贡献。
% 与第 1 章"研究意义"呼应，但这里用完成时态。

\subsection{研究局限}

% 诚实说明研究的不足之处。
% 例如：样本范围有限、方法存在假设约束、数据时间跨度不够等。

\subsection{未来展望}

% 基于局限性提出后续研究方向。
% 可以指出值得深入探索的问题或可以改进的方法。


\newpage
\addcontentsline{toc}{section}{参考文献}
\defbibheading{bibliography}[\bibname]{\section*{#1}\bfseries\zihao{4}\heiti\normalfont}
\printbibliography[heading=bibliography,title=参考文献]

\newpage
\addcontentsline{toc}{section}{致谢}
\begin{center}
	\zihao{4}\bfseries{致谢}
\end{center}
\vspace{\baselineskip}

\zihao{-4}{这里填写致谢内容。}

\end{document}
